% solutions/a3-product-hierarchical-prior.tex
% Standalone proof dossier for two closely coupled A3 nodes.
\documentclass[../main.tex]{subfiles}
\begin{document}

% === SOLUTION HEADER (proof provenance) =====================================
%   ledger-node : thm:a3-product; prop:a3-hierarchical-prior
%   refines     : thm:a3-product; prop:a3-hierarchical-prior
%   bounded_by  : obs:heavy-tail-no-classical; obs:marginals-not-joint (both nodes)
%   checked_by  : agent
%   author      : /root/a3_prior_prover
%   reviewer    : /root/a3_prior_review
%   review      : research/reviews/2026-08-27-a3-product-hierarchical-prior-post-promotion-proof-review.md
%   date        : 2026-08-27
% ============================================================================

\section*{Solution: product Hardy tensorization and the hierarchical-prior pullback}
\label{sec:sol-a3-product-hierarchical-prior}

\paragraph{Refined statements.}
The first theorem proves the product Hardy certificate of
Theorem~\ref{thm:a3-product}, including the exact tensorization rule before the Hardy
factor-four estimate is applied.  The second proves the prior-side statement of
Proposition~\ref{prop:a3-hierarchical-prior}.  Throughout, all products are finite.

\begin{theorem}[Product Hardy certificate]
\label{thm:sol-a3-product}
Let $n\ge1$, and for each $i\le n$ let $\mu_i(\dd x)=p_i(x)\dd x$ be a probability law on
$\R$ with median $m_i$, where the expressions below are understood in the extended sense if
$p_i$ vanishes.  Let $a_i:\R\to(0,\infty)$ be a measurable weight and set
\begin{align*}
 B_{i,+}
 &=\sup_{x>m_i}\mu_i([x,\infty))
             \int_{m_i}^x\frac{\dd t}{a_i(t)p_i(t)},\\
 B_{i,-}
 &=\sup_{x<m_i}\mu_i(({-\infty},x])
             \int_x^{m_i}\frac{\dd t}{a_i(t)p_i(t)},\\
 B_i&=\max(B_{i,+},B_{i,-}).
\end{align*}
If $B_i<\infty$ for every $i$, then $\mu=\bigotimes_{i=1}^n\mu_i$ satisfies
\begin{equation}
\label{eq:sol-a3-product-hardy}
 \Var_\mu(f)
 \le 4\Big(\max_{1\le i\le n}B_i\Big)
 \int_{\R^n}\sum_{i=1}^n a_i(x_i)\abs{\partial_i f(x)}^2\dd\mu(x).
\end{equation}
More generally, if the $i$th factor has optimal weighted Poincar\'e constant $C_i$ for
the form $\int a_i\abs{g'}^2\dd\mu_i$, then the optimal constant of the product in the
displayed diagonal metric is exactly $\max_i C_i$.
\end{theorem}

\begin{proof}
For a smooth $f$ in the product form domain, the Efron--Stein (conditional-variance)
inequality for independent coordinates gives
\begin{equation}
\label{eq:sol-a3-efron-stein}
 \Var_\mu(f)
 \le \sum_{i=1}^n
 \int_{\R^{n-1}}
 \Var_{\mu_i}\!\big(f(x_{-i},\mathord\cdot)\big)\dd\mu_{-i}(x_{-i}).
\end{equation}
For completeness, the two-factor case follows from
\[
 \Var_{\mu_1\otimes\mu_2}(f)
 =\int\Var_{\mu_1}(f(\mathord\cdot,x_2))\dd\mu_2(x_2)
  +\Var_{\mu_2}\!\left(\int f(x_1,\mathord\cdot)\dd\mu_1(x_1)\right).
\]
Convexity of variance, or conditional Jensen, bounds the last term by
$\int\Var_{\mu_2}(f(x_1,\mathord\cdot))\dd\mu_1(x_1)$.  Iterating this two-factor
decomposition proves \eqref{eq:sol-a3-efron-stein} for every finite product.

If the $i$th coordinate inequality holds with constant $C_i$, applying it inside the $i$th
term of \eqref{eq:sol-a3-efron-stein} and then using Fubini yields
\[
 \Var_\mu(f)
 \le \sum_{i=1}^n C_i\int a_i(x_i)\abs{\partial_i f}^2\dd\mu
 \le \big(\max_iC_i\big)
       \int\sum_{i=1}^n a_i(x_i)\abs{\partial_i f}^2\dd\mu.
\]
Conversely, testing the product inequality on $f(x)=g(x_k)$ shows that its optimal constant is
at least the optimal constant $C_k$ of every factor.  Approximate minimizers suffice if a
factor gap is not attained.  Thus the product constant is exactly $\max_iC_i$.

The weighted one-dimensional Hardy--Muckenhoupt theorem underlying
Theorem~\ref{thm:hardy-1d} gives $C_i\le4B_i$ for every factor.  Substitution in the exact
tensorization rule proves \eqref{eq:sol-a3-product-hardy}.  The argument starts on a smooth core
and extends to the closed product form by approximation.
\end{proof}

\begin{proposition}[Hierarchical horseshoe prior in the pullback metric]
\label{prop:sol-a3-hierarchical-prior}
Fix $d\ge1$.  Let $Z_1,\ldots,Z_d$ be independent $N(0,1)$ variables.  Independently, let
$S_0,S_1,\ldots,S_d$ be standard half-Cauchy variables, and write
\[
 U=\log S_0,\qquad V_j=\log S_j,\qquad
 \Theta_j=e^{U+V_j}Z_j.
\]
Let $\pi$ be the law of $(\Theta,U,V)$ on
$\R^d\times\R\times\R^d$.  Its optimal weighted Poincar\'e constant is $4$ for the
pullback energy
\begin{equation}
\label{eq:sol-a3-pullback-energy}
 \Gamma_{\mathrm{pb}}(f)
 =\sum_{j=1}^d e^{2u+2v_j}\abs{\partial_{\theta_j}f}^2
  +\sum_{j=1}^d
    \abs{\partial_{v_j}f+\theta_j\partial_{\theta_j}f}^2
  +\abs{\partial_uf+\sum_{j=1}^d
                  \theta_j\partial_{\theta_j}f}^2.
\end{equation}
That is,
\begin{equation}
\label{eq:sol-a3-hierarchical-pi}
 \Var_\pi(f)\le4\int\Gamma_{\mathrm{pb}}(f)\dd\pi,
\end{equation}
and no smaller constant works on the associated closed form domain.
\end{proposition}

\begin{proof}
We first identify and normalize the log-scale law.  A standard half-Cauchy variable $S$ has
density $2/\{\pi(1+s^2)\}$ on $(0,\infty)$.  Therefore $Y=\log S$ has density
\begin{equation}
\label{eq:sol-a3-log-half-cauchy}
 q(y)=\frac{2e^y}{\pi(1+e^{2y})}
     =\frac1{\pi\cosh y},\qquad y\in\R.
\end{equation}
The change of variables already proves that $q$ integrates to one (equivalently,
$\int_\R\operatorname{sech}y\,\dd y=\pi$).

Let $X$ instead be a standard Cauchy variable and put $Y_*=\operatorname{arsinh}X$.
Since $x=\sinh y$ has Jacobian $\cosh y$,
\[
 \frac1{\pi(1+\sinh^2y)}\cosh y
 =\frac1{\pi\cosh y}=q(y).
\]
Thus $\log S$ and $\operatorname{arsinh}X$ are equal in law.  They are not being identified
pathwise; the equality of their normalized pushforward densities is what is used below.

Apply Theorem~\ref{thm:a3-student} with $d=1$ and $\beta=1$.  Its sharp spectral gap is
$\lambda_{1,1}=(1-\tfrac12)^2=\tfrac14$, so the standard Cauchy law $\mathsf C$ obeys the sharp
inequality
\begin{equation}
\label{eq:sol-a3-cauchy-sharp}
 \Var_{\mathsf C}(g)
 \le4\int_\R(1+x^2)\abs{g'(x)}^2\dd\mathsf C(x).
\end{equation}
For $g(x)=h(\operatorname{arsinh}x)$,
$g'(x)=h'(\operatorname{arsinh}x)/\sqrt{1+x^2}$.  Hence the change of variables just computed
turns \eqref{eq:sol-a3-cauchy-sharp} into
\begin{equation}
\label{eq:sol-a3-sech-sharp}
 \Var_q(h)\le4\int_\R\abs{h'(y)}^2q(y)\dd y.
\end{equation}
The correspondence $h(y)\leftrightarrow g(x)=h(\operatorname{arsinh}x)$ is an isometry of the
two Dirichlet form domains, in both $L^2$ norm and energy.  Therefore sharpness in
\eqref{eq:sol-a3-cauchy-sharp} also proves that the optimal constant in
\eqref{eq:sol-a3-sech-sharp} is exactly $4$.

We also recall the elementary standard-Gaussian Poincar\'e inequality with constant $1$.
One short proof uses the Ornstein--Uhlenbeck semigroup
$P_tg(x)=\E[g(e^{-t}x+\sqrt{1-e^{-2t}}G)]$.  Invariance of the standard Gaussian law $\gamma$,
the identity $(P_tg)'=e^{-t}P_t(g')$, and variance dissipation give
\[
 \Var_\gamma(g)
 =2\int_0^\infty\!\int\abs{(P_tg)'(x)}^2\dd\gamma(x)\dd t
 \le2\int_0^\infty e^{-2t}\dd t\int\abs{g'}^2\dd\gamma
 =\int\abs{g'}^2\dd\gamma.
\]
(The coordinate function, by form-domain approximation, shows sharpness.)

Let
\[
 \nu=\gamma^{\otimes d}\otimes q\otimes q^{\otimes d}
\]
be the product law of $(Z,U,V)$.  Applying the same conditional-variance tensorization as in
\eqref{eq:sol-a3-efron-stein} to the Gaussian constant $1$ and the log-scale constant $4$
shows that $\nu$ has Euclidean Poincar\'e constant at most $4$.  It is exactly $4$: by
sharpness of \eqref{eq:sol-a3-sech-sharp}, there is a sequence $h_k$ in the $q$-form domain
whose variance-to-energy ratios tend to $4$, and the functions $(z,u,v)\mapsto h_k(u)$ have the
same ratios under $\nu$.  No assertion that a single extremizer exists is needed.

Now consider the smooth bijection
\[
 T(z,u,v)=(\theta,u,v),\qquad \theta_j=e^{u+v_j}z_j.
\]
Its inverse is $z_j=e^{-u-v_j}\theta_j$, and by construction $T_\#\nu=\pi$.  For a smooth
$f(\theta,u,v)$ put $F=f\circ T$.  Derivatives of $f$ on the right below are evaluated at
$T(z,u,v)$, with $\theta$ held fixed in the explicit $u$- and $v$-derivatives.  The complete
chain rule is
\begin{equation}
\label{eq:sol-a3-chain-rule}
 \partial_{z_j}F=e^{u+v_j}\partial_{\theta_j}f,
 \qquad
 \partial_{v_j}F=\partial_{v_j}f+\theta_j\partial_{\theta_j}f,
 \qquad
 \partial_uF=\partial_uf+\sum_{j=1}^d\theta_j\partial_{\theta_j}f.
\end{equation}
Consequently $\abs{\nabla_{z,u,v}F}^2=\Gamma_{\mathrm{pb}}(f)\circ T$.  To make all structural
cross terms explicit, the same quantity expands as
\begin{align}
\label{eq:sol-a3-expanded-cross-terms}
 \Gamma_{\mathrm{pb}}(f)
 ={}&\sum_{j=1}^d(e^{2u+2v_j}+2\theta_j^2)
             \abs{\partial_{\theta_j}f}^2
     +\sum_{j=1}^d\abs{\partial_{v_j}f}^2+\abs{\partial_uf}^2 \\
 &+2\sum_{j=1}^d\theta_j(\partial_{v_j}f)(\partial_{\theta_j}f)
   +2\sum_{j=1}^d\theta_j(\partial_uf)(\partial_{\theta_j}f)
   +2\sum_{1\le j<k\le d}\theta_j\theta_k
       (\partial_{\theta_j}f)(\partial_{\theta_k}f).
\end{align}
In particular, the coefficient--local-scale, coefficient--global-scale, and distinct-coefficient
cross terms are all retained.

The Euclidean inequality for $\nu$, followed by the pushforward identity, now gives
\[
 \Var_\pi(f)=\Var_\nu(F)
 \le4\int\abs{\nabla F}^2\dd\nu
 =4\int\Gamma_{\mathrm{pb}}(f)\dd\pi,
\]
which proves the upper bound.  The diffeomorphism carries smooth cores into smooth cores, and
the identity of energies extends this inequality to the closed pullback form domain.

Finally, use the same near-extremizing sequence as above and set
$f_k(\theta,u,v)=h_k(u)$.  Then
$\Gamma_{\mathrm{pb}}(f_k)=\abs{h_k'(u)}^2$, while the $u$-marginal of $\pi$ is $q$.
The corresponding variance-to-energy ratios tend to $4$.  Thus no constant smaller than $4$
can satisfy \eqref{eq:sol-a3-hierarchical-pi}, completing the optimality proof.
\end{proof}

\paragraph{Obstructions respected.}
Both \texttt{thm:a3-product} and \texttt{prop:a3-hierarchical-prior} have the formal
\texttt{bounded\_by} edges \texttt{obs:heavy-tail-no-classical} and
\texttt{obs:marginals-not-joint}.  The statements respect those fences, and the neighboring
\texttt{obs:flat-direction} context, as follows:
\begin{itemize}[leftmargin=2em]
\item \texttt{obs:heavy-tail-no-classical}: Theorem~\ref{thm:sol-a3-product} uses a weighted diagonal
  metric, and Proposition~\ref{prop:sol-a3-hierarchical-prior} uses Euclidean geometry only in
  Gaussian/log-scale coordinates, equivalently the non-Euclidean pullback form
  \eqref{eq:sol-a3-pullback-energy}.  Neither asserts a raw Euclidean gap for half-Cauchy scales.
\item \texttt{obs:marginals-not-joint}: the product theorem assumes independence.  The hierarchical
  result treats one fully specified dependent pushforward of an independent base law and retains
  every scale derivative and pullback cross term; it makes no marginal-only claim for arbitrary
  dependence.
\item \texttt{obs:flat-direction}: the proposition is a prior statement with no likelihood.  It makes no
  assertion that a saturating or weakly identifying likelihood supplies a posterior curvature
  floor.
\end{itemize}

\paragraph{Certification boundary.}
The two results are unconditional relative to the imported theorems recorded in their ledger
closures: the weighted Hardy--Muckenhoupt criterion (\texttt{thm:hardy-1d}) and the sharp
generalized-Cauchy theorem (\texttt{thm:a3-student}).  No numerical evidence is used.
This dossier is independently agent-certified for exactly the two ledger nodes named in its
header; it makes no posterior-likelihood or arbitrary-dependence claim.

\end{document}
